% Tóm tắt (tiếng Việt, một trang): nội dung tương ứng Abstract.
\FrontChapter{TÓM TẮT}
\begin{otherlanguage}{vietnamese}
Degron là đoạn trình tự đánh dấu protein để bị phân hủy có chọn lọc; một phép thay thế amino acid trong
degron có thể làm thay đổi hoạt tính này. Thực nghiệm thường chấm điểm cả peptide, trong khi đo mọi
phép thay thế rất tốn kém. Luận văn nghiên cứu cách dự đoán hiệu ứng của phép thay thế đơn amino acid
trong peptide degron ở mức residue từ dữ liệu độ bền protein thực nghiệm, đánh giá khả năng tổng quát hóa
khi chia dữ liệu theo cụm protein so với các đối chứng ridge, chênh lệch điểm PAP và Degformer, đồng
thời xây dựng nguyên mẫu hiển thị ma trận hiệu ứng, xếp hạng phép thay thế và xuất kết quả kèm nguồn gốc.
\textbf{[CHỜ BỔ SUNG: kết quả chính và giới hạn, viết khi hoàn thành công việc.]}

\par\vspace{6bp}\noindent\textbf{Từ khóa:} degron, độ bền protein, dự đoán hiệu ứng đột biến, hệ ubiquitin--proteasome, học máy.
\end{otherlanguage}
